\section*{Abstract}
Video-based learning has transformed education in Bangladesh, yet millions of technical lectures remain inaccessible as searchable, reviewable text due to the unique challenges of Banglish---the pervasive code-mixing of Bengali and English in educational discourse. This thesis presents a Multimodal Banglish Classroom Summarizer, an end-to-end system that transforms lecture videos into structured notes by fusing automatic speech recognition with visual content extraction. The system employs OpenAI's Whisper for speech transcription, Qwen2.5-VL-7B for whiteboard keyword extraction, and cross-modal verification to correct ASR errors using visual ground truth. We introduce the BanglaASR benchmark dataset comprising nine technical lectures (73,141 characters) across Python programming, Digital Logic Design, and Database Management Systems, along with term-based evaluation metrics that accommodate the orthographic variability inherent to Romanized Banglish. Experimental evaluation demonstrates 73.9\% Term F1, 83.6\% Term Precision, and 66.5\% Term Recall, with cross-modal fusion improving baseline Whisper performance by 5.7 percentage points. Our analysis reveals that verification-oriented fusion outperforms naive visual biasing, which introduces hallucinations that degrade accuracy. The system operates on consumer-grade GPU hardware (NVIDIA RTX 3090, 24GB VRAM) through careful memory management. This work establishes reproducible baselines for Banglish technical lecture processing and contributes both methodological insights and practical tools for an underserved linguistic community of over 230 million Bengali speakers.

\vspace{1cm}
\textbf{Keywords:} Automatic Speech Recognition, Code-Mixed Language Processing, Banglish, Multimodal Learning, Vision-Language Models, Lecture Transcription, Educational Technology
\pagebreak
